\begin{figure}[htbp]
\centering
\begin{tikzpicture}[>=Latex,font=\small]
 \node at (1.4,4.75) {reference element $\widehat K$};
 \node at (7.4,4.75) {curved physical element $K$};
 \foreach \s in {0,.25,.5,.75,1} {
  \draw[blue!30] ({2.8*\s},0)--({2.8*\s},2.8);
  \draw[blue!30] (0,{2.8*\s})--(2.8,{2.8*\s});
  \draw[blue!30,domain=0:1,samples=35]
    plot ({6+2.8*\x},{2.8*((1+.25*\x)*\s+.2*\x*\x)});
  \draw[blue!30] ({6+2.8*\s},{2.8*.2*\s*\s})--
    ({6+2.8*\s},{2.8*(1+.25*\s+.2*\s*\s)});
 }
 \draw[blue!70!black,thick] (0,0) rectangle (2.8,2.8);
 \foreach \s in {0,1} {
  \draw[blue!70!black,thick,domain=0:1,samples=35]
    plot ({6+2.8*\x},{2.8*((1+.25*\x)*\s+.2*\x*\x)});
  \draw[blue!70!black,thick] ({6+2.8*\s},{2.8*.2*\s*\s})--
    ({6+2.8*\s},{2.8*(1+.25*\s+.2*\s*\s)});
 }
 \foreach \s in {.2,.5,.8} {
  \draw[->,orange!85!black,thick]
    ({6+2.8*\s},{2.8*(1+.25*\s+.2*\s*\s)})--
    ({6+2.8*\s-.52*(.25+.4*\s)},{2.8*(1+.25*\s+.2*\s*\s)+.52});
 }
 \draw[->,thick] (3.25,1.55)--(5.4,1.55);
 \node[above] at (4.3,1.55) {$T_K$};
 \node[below,align=center] at (1.4,-.25) {straight reference grid\\coordinates $(\xi,\eta)$};
 \node[below,align=center] at (7.4,-.25) {mapped grid\\spatially varying Jacobian and normals};
\end{tikzpicture}
\caption{An illustrative polynomial map
$T_K(\xi,\eta)=(\xi,(1+\xi/4)\eta+\xi^2/5)$ from the unit square.
Its determinant is $1+\xi/4>0$: equal reference areas need not have equal
physical areas. Orange arrows show outward normal directions along the curved
top face (not their magnitudes). The interior lines are coordinate lines
within one element, not separate DG cells. This is a geometric schematic,
not a mesh from a numerical experiment.}
\label{fig:curved-map}
\end{figure}
